% part: intuitionistic-logic
% chapter: propositions-as-types
% section: sequent-natural-deduction

\documentclass[../../../include/open-logic-section]{subfiles}

\begin{document}

\olfileid[hi]{int}{pty}{snd}

\olsection{अनुक्रम प्राकृतिक निगमन}

यदि कोई ऐसी प्राकृतिक निगमन !!{derivation} है जिसमें $\Gamma$ मुक्त न की गई
मान्यताएँ और $!A$ निष्कर्ष है, तो $\Gamma \Sequent !A$ लिखें; और यदि $\Gamma$
रिक्त है तो $\Sequent !A$।

$\Gamma \cup \{!A_1, \dots, !A_n\}$ के लिए $\Gamma, !A_1, \dots, !A_n$,
और $\Gamma \cup \Delta$ के लिए $\Gamma, \Delta$ लिखते हैं।

ध्यान दें कि $\Gamma \Sequent !A \land !A$ होने का अर्थ है कि हमारे पास ऐसी
!!a{derivation} है जिसमें $\Gamma$ मुक्त न की गई मान्यताएँ और $!A \land !A$
अंतिम !!{formula} है। तब नीचे $\Elim{\land}$ लगाकर हम समान मुक्त न की गई
मान्यताओं और निष्कर्ष $!A$ वाली !!a{derivation} पा सकते हैं। दूसरे शब्दों में,
यदि $\Gamma \Sequent !A \land !B$, तो $\Gamma \Sequent !A$।
\begin{prooftree}
  \Axiom$\Gamma \fCenter !A \land !B$
  \RightLabel{$\Elim{\land}$}
  \UnaryInf$\Gamma \fCenter !A$
  \DisplayProof\qquad\bottomAlignProof
  \Axiom$\Gamma \fCenter !A \land !B$
  \RightLabel{$\Elim{\land}$}
  \UnaryInf$\Gamma \fCenter !B$
\end{prooftree}
$\Elim{\land}$ चिह्न प्राकृतिक निगमन के इसी नाम के नियम से संबंध का संकेत
देता है।

इसी प्रकार मान लें $\Gamma, !A \Sequent !B$, अर्थात् हमारे पास मुक्त न की गई
मान्यताओं $\Gamma, !A$ और अंतिम !!{formula}~$!B$ वाली !!a{derivation} है। यदि
हम $\Intro{\lif}$ नियम लगाएँ, तो हमें ऐसी !!a{derivation} मिलती है जिसमें
$\Gamma$ मुक्त न की गई मान्यताएँ और $!A \lif !B$ अंतिम !!{formula} है,
अर्थात् $\Gamma \Sequent !A \lif !B$। ध्यान दें कि इससे मान्यताओं का
!!{discharge} अधिक स्पष्ट हो गया है।
\begin{prooftree}
  \Axiom $\Gamma, !A \fCenter !B$
  \RightLabel{$\Intro{\lif}$}
  \UnaryInf $\Gamma \fCenter !A \lif !B$
\end{prooftree}

अन्य नियमों से भी इसी प्रकार निष्कर्ष निकाल सकते हैं, जिन्हें नीचे स्पष्ट किया
गया है:
\begin{gather*}
  \Axiom $\Gamma \fCenter !A$
  \Axiom $\Delta \fCenter !B$
  \RightLabel{$\Intro{\land}$}
  \BinaryInf $\Gamma,\Delta \fCenter !A \land !B$
  \DisplayProof\\
  \Axiom $\Gamma \fCenter !A \land !B$
  \RightLabel{$\Elim{\land}_1$}
  \UnaryInf $\Gamma \fCenter !A$
  \DisplayProof
  \qquad
  \Axiom $\Gamma \fCenter !A \land !B$
  \RightLabel{$\Elim{\land}_2$}
  \UnaryInf $\Gamma \fCenter !B$
  \DisplayProof
  \\
  \Axiom $\Gamma \fCenter !A$
  \RightLabel{$\Intro{\lor}_1$}
  \UnaryInf $\Gamma \fCenter !A \lor !B$
  \DisplayProof\qquad
  \Axiom $\Gamma \fCenter !B$
  \RightLabel{$\Intro{\lor}_2$}
  \UnaryInf $\Gamma \fCenter !A \lor !B$
  \DisplayProof\\
  \Axiom $\Gamma \fCenter !A \lor !B$
  \Axiom $\Delta, !A \fCenter !C$
  \Axiom $\Delta', !B \fCenter !C$
  \RightLabel{$\Elim{\lor}$}
  \TrinaryInf $\Gamma, \Delta, \Delta' \fCenter !C$
  \DisplayProof
  \\
  \Axiom $\Gamma, !A \fCenter !B$
  \RightLabel{$\Intro{\lif}$}
  \UnaryInf $\Gamma \fCenter !A \lif !B$
  \DisplayProof
  \qquad
  \Axiom $\Delta \fCenter !A \lif !B$
  \Axiom $\Gamma \fCenter !A$
  \RightLabel{$\Elim{\lif}$}
  \BinaryInf $\Gamma, \Delta \fCenter !B$
  \DisplayProof
  \\
  \Axiom $\Gamma \fCenter \lfalse$
  \RightLabel{$\FalseInt$}
  \UnaryInf $\Gamma \fCenter !A$
  \DisplayProof
\end{gather*}

हर मान्यता अपने आप~$!A$ से $!A$ की !!a{derivation} है, अर्थात् हमारे पास सदा
$!A \Sequent !A$ है।

\begin{prooftree}
  \AxiomC{$ $}
  \UnaryInfC{$!A \fCenter !A$}
\end{prooftree}

इन नियमों को मिलाकर इस विषय का कलन माना जा सकता है कि कौन-सी प्राकृतिक निगमन
!!{derivation}s विद्यमान हैं। इन्हें प्राकृतिक निगमन का संकेतनात्मक रूपांतर भी
माना जा सकता है, जिसमें प्रत्येक चरण न केवल !!{derive} किए गए !!{formula} को,
बल्कि उन मुक्त न की गई मान्यताओं को भी दर्ज करता है जिनसे उसे !!{derive} किया
गया।

\begin{prooftree}
  \Axiom$!A \fCenter !A$
  \UnaryInf $!A \fCenter !A \lor (!A \lif \lfalse)$
  \Axiom $!B \fCenter !B$
  \BinaryInf$!A, !B\lif \fCenter \lfalse$
  \UnaryInf$(!B \fCenter !A \lif \lfalse$
  \UnaryInf$( !B \fCenter !A \lor (!A \lif \lfalse)$
  \Axiom$(!B \fCenter !B$
  \BinaryInf$(!B \fCenter \lfalse$
  \UnaryInf$\fCenter !B \lif \lfalse$
\end{prooftree}
जहाँ $!B$, $(!A \lor (!A \lif \lfalse))\lif \lfalse$ का संक्षिप्त रूप है।

\end{document}
